\begin{corollary}[The actual posterior mean can have unbounded range]
Use the four observations in the preceding construction, so $n=4$.
There are absolute constants $c,c',\delta_0>0$ such that, for
$0<\delta<\delta_0$, $\lambda=\delta^2$, every $N\ge4$, and every
$0<\nu\le1$, the ordinary Gaussian posterior mean satisfies
\[
 \max\{\bar z(-r+\ell),\bar z(r-\ell)\}\ge c/\delta. \tag{RG4}
\]
Here $r$ and $\ell$ are the fixed constants in the preceding proof.
Moreover, against the uniform test-input law on $[-1,1]$ and any test
labels bounded by one, its squared test risk is at least
$c'/\delta^2$. The posterior-average squared test risk has the same
lower bound.

This is a deterministic-training-sample statement. It is not a lower
bound with a stated probability under an IID sample from that test law,
and it is not an impossibility theorem for other likelihood schedules.
\end{corollary}
\begin{proof}
Write $K=r+\sqrt3$. The explicit interpolant has zero empirical loss
and path variation at most $K/\delta$. Proposition (PR2) gives
\[
 \mathbb E_\Pi\widehat R(z)+2\lambda\mathbb E_\Pi P
 \le2\lambda K/\delta+C\lambda\nu
             \log\frac e{\lambda\nu}.
\]
For $\lambda=\delta^2$ and $0<\nu\le1$, this implies, uniformly in
both width and scale,
\[
 \mathbb E_\Pi\widehat R(z)\le C\delta,
 \qquad
 \mathbb E_\Pi P\le K/\delta+C\log(e/\delta). \tag{RG5}
\]
Let $f_i=z(X_i)$ in the order of the four training observations, and
let $f_L=z(-r+\ell)$, $f_R=z(r-\ell)$. The triangular-hat identity,
without assuming exact interpolation, is
\[
 \frac1a\int(\psi_L+\psi_R)\,dz''
 =\frac{f_1+f_4-f_2-f_3}{\delta}
   +\frac{f_L+f_R-f_2-f_3}{\ell}.
\]
Since the negative of this integral is at most $P$, rearrangement
and posterior expectation give
\begin{align*}
 \bar z(-r+\ell)+\bar z(r-\ell)
 &\ge \frac{2\ell}{\delta}
       -\frac{\ell}{a}\mathbb E_\Pi P
       -C(1+\delta^{-1})
                       \sqrt{\mathbb E_\Pi\widehat R(z)}\\
 &\ge\frac{\ell(2-K/a)}{\delta}
       -C\delta^{-1/2}-C\log(e/\delta).
\end{align*}
The leading coefficient is strictly positive by the choice of
$\ell$. For sufficiently small $\delta$, this proves (RG4).
All expectations are finite because the Gaussian posterior density
is bounded by $Z^{-1}$ times its prior density.

Every normalized ridge is $1$-Lipschitz, so the posterior mean is
Lipschitz with constant at most $\mathbb E_\Pi P\le2K/\delta$
for small $\delta$. A value at least $c/\delta$ at either fixed
interior point therefore persists at height $c/(2\delta)$ on an
interval of fixed positive length. For bounded test labels and
small enough $\delta$, integrating the squared error on that interval
gives $c'/\delta^2$. Jensen's inequality gives the posterior-average
risk lower bound as well.
\end{proof}
